\documentclass[a4paper,12pt]{article}
\usepackage[utf8]{inputenc}
\usepackage[T2A]{fontenc}
\usepackage[russian]{babel}
\usepackage{amsmath,amssymb}
\usepackage{geometry}
\usepackage{hyperref}
\usepackage{enumitem}
\usepackage{parskip}
\geometry{margin=2.5cm}

\title{Обработка изображений и видео\\\large Учебное пособие по лабораторной работе}
\author{}
\date{}

\begin{document}
\maketitle

\section*{Введение}

Цифровое изображение --- это двумерный массив чисел. Каждое число описывает яркость или цвет точки изображения в определённой системе координат. Работа с изображением сводится к математическим операциям над этим массивом: изменению значений, вычислению новых величин на основе соседних пикселей, поиску характерных областей. Видео --- это последовательность кадров, поэтому все операции над отдельным изображением применимы и к видео, но добавляется временнáя ось: кадры сравниваются между собой для выделения движения.

Настоящее пособие описывает математическую основу всех операций лабораторной работы: от базового представления цвета до алгоритмов выделения ключевых точек и вычитания фона на видео.

\section{Представление цвета}

\subsection{Аддитивная модель RGB}

Модель RGB основана на принципе аддитивного смешения трёх базовых цветов: красного (Red), зелёного (Green), синего (Blue). Любой цвет представляется как линейная комбинация этих трёх компонент:
\[
C = R \cdot \mathbf{e}_R + G \cdot \mathbf{e}_G + B \cdot \mathbf{e}_B,
\]
где $R, G, B \in [0, 255]$ для 8-битного представления. Значение $0$ соответствует полному отсутствию компоненты, $255$ --- её максимальной интенсивности. Белый цвет получается при $R = G = B = 255$, чёрный --- при $R = G = B = 0$. Модель RGB аддитивна, потому что смешение света разных длин волн даёт суммарный световой поток; именно так работают мониторы и проекторы.

\subsection{Субтрактивная модель CMYK}

Модель CMYK описывает не излучаемый, а отражённый свет. Голубая (Cyan), пурпурная (Magenta), жёлтая (Yellow) краски поглощают (вычитают) соответствующие длины волн из падающего белого света, а чёрная (Key) добавляется отдельно для получения глубоких теней без использования смеси трёх красок. Переход из RGB в CMYK:
\[
K = 1 - \max(R', G', B'), \quad
C = \frac{1 - R' - K}{1 - K}, \quad
M = \frac{1 - G' - K}{1 - K}, \quad
Y = \frac{1 - B' - K}{1 - K},
\]
где $R', G', B'$ --- нормированные к единице значения RGB. Если $K = 1$, все остальные компоненты полагаются равными нулю (полностью чёрный цвет). Обратное преобразование:
\[
R = 255(1 - C)(1 - K), \quad
G = 255(1 - M)(1 - K), \quad
B = 255(1 - Y)(1 - K).
\]

\subsection{Модели HSL и HSV}

В моделях HSL (Hue, Saturation, Lightness) и HSV (Hue, Saturation, Value) цвет описывается не через смесь базовых компонент, а через положение на цветовом круге и два параметра, отвечающих за насыщенность и светлоту. Это удобнее для задач, где нужно изменить один аспект цвета, не затрагивая остальные.

Оттенок $H$ задаёт угол на цветовом круге от $0^\circ$ до $360^\circ$: $0^\circ$ --- красный, $120^\circ$ --- зелёный, $240^\circ$ --- синий. Насыщенность $S \in [0, 100]$ показывает, насколько цвет далёк от серого: $S = 0$ соответствует полному отсутствию цветности, $S = 100$ --- максимально чистому цвету.

В модели HSL светлота $L$ определяется как среднее между максимумом и минимумом компонент:
\[
L = \frac{\max(R', G', B') + \min(R', G', B')}{2}.
\]
Насыщенность вычисляется по формуле
\[
S = \begin{cases}
\dfrac{\max - \min}{\max + \min}, & L \le 0{,}5, \\[6pt]
\dfrac{\max - \min}{2 - \max - \min}, & L > 0{,}5.
\end{cases}
\]
Оттенок $H$ вычисляется по-разному в зависимости от того, какая из компонент максимальна.

Модель HSV отличается только третьей компонентой: вместо светлоты используется \emph{значение} $V = \max(R', G', B')$. Модель HSV удобна для алгоритмов цветовой сегментации, HSL --- для изменения яркости в интерфейсах.

\subsection{Модель LAB}

Пространство CIE L*a*b* разработано для того, чтобы евклидово расстояние между двумя цветами соответствовало воспринимаемому различию. Компонента $L^*$ описывает светлоту от $0$ (чёрный) до $100$ (белый). Компоненты $a^*$ и $b^*$ описывают цветность: положительные значения $a^*$ смещают цвет к пурпурному, отрицательные --- к зелёному; положительные $b^*$ --- к жёлтому, отрицательные --- к синему.

Переход из RGB в LAB выполняется через промежуточное пространство CIE XYZ:
\[
\begin{pmatrix} X \\ Y \\ Z \end{pmatrix}
= M_{sRGB} \begin{pmatrix} R' \\ G' \\ B' \end{pmatrix},
\]
где $R', G', B'$ --- значения, к которым применена гамма-коррекция (см.~раздел о профилях), а матрица $M_{sRGB}$ --- матрица перехода, определённая стандартом sRGB. Далее нормированные координаты переводятся к виду
\[
f(t) = \begin{cases} t^{1/3}, & t > 0{,}008856, \\ 7{,}787 t + 16/116, & t \le 0{,}008856. \end{cases}
\]
Окончательные компоненты:
\[
L^* = 116 f(Y / Y_n) - 16, \quad
a^* = 500 [f(X / X_n) - f(Y / Y_n)], \quad
b^* = 200 [f(Y / Y_n) - f(Z / Z_n)],
\]
где $X_n, Y_n, Z_n$ --- координаты белой точки (для sRGB --- $D_{65}$).

\subsection{Модель YCbCr}

Модель YCbCr разделяет яркость и цветность. Это удобно для сжатия, поскольку человеческий глаз чувствителен к яркости, но плохо различает мелкие изменения цветности. Компонента $Y$ вычисляется как взвешенная сумма каналов RGB:
\[
Y = 0{,}299 R + 0{,}587 G + 0{,}114 B.
\]
Коэффициенты подобраны согласно стандарту ITU-R BT.601 с учётом чувствительности глаза к разным длинам волн. Цветоразностные компоненты:
\[
Cb = 128 - 0{,}168736 R - 0{,}331264 G + 0{,}5 B, \quad
Cr = 128 + 0{,}5 R - 0{,}418688 G - 0{,}081312 B.
\]
Смещение на $128$ нужно, чтобы цветоразностные компоненты могли принимать отрицательные значения в диапазоне $[0, 255]$. Обратное преобразование:
\[
R = Y + 1{,}402 (Cr - 128), \quad
G = Y - 0{,}344136 (Cb - 128) - 0{,}714136 (Cr - 128), \quad
B = Y + 1{,}772 (Cb - 128).
\]

\section{Цветовые профили}

Цветовой профиль описывает, как числа RGB связаны с реальным физическим цветом для конкретного устройства. Одно и то же значение $(200, 100, 50)$ на разных мониторах отображается по-разному; профиль задаёт словарь перевода между кодами устройства и абсолютными координатами CIE XYZ.

Профиль определяется несколькими параметрами. \emph{Первичные цвета} задают координаты вершин цветового охвата. \emph{Белая точка} указывает координаты цвета, который устройство считает чистым белым. \emph{Гамма-кривая} описывает нелинейность отклика.

Для стандарта sRGB гамма-кривая имеет кусочно-заданный вид:
\[
\gamma_{sRGB}(v) = \begin{cases}
v / 12{,}92, & v \le 0{,}04045, \\
\left(\dfrac{v + 0{,}055}{1{,}055}\right)^{2{,}4}, & v > 0{,}04045.
\end{cases}
\]
Нелинейность нужна потому, что человеческое восприятие яркости логарифмично, а физическая яркость линейна. Без гамма-коррекции тёмные участки изображения выглядели бы слишком тёмными.

Разные профили описываются разными матрицами перехода RGB $\to$ XYZ. Например, для Adobe RGB матрица отличается от sRGB, поэтому коэффициенты в формуле перехода к LAB тоже будут другими. Это и есть связь между профилем и коэффициентами: коэффициенты --- это и есть профиль, записанный в виде чисел.

\section{Преобразование в градации серого}

Преобразование в градации серого сводит трёхканальное изображение к одноканальному. Каждый пиксель $(R, G, B)$ заменяется на значение яркости
\[
Y = 0{,}299 R + 0{,}587 G + 0{,}114 B.
\]
Коэффициенты подобраны так, что сумма равна единице, поэтому белый цвет $(255, 255, 255)$ превращается в $255$, а чёрный --- в $0$. Вес зелёного канала наибольший, потому что человеческий глаз наиболее чувствителен к средним длинам волн; вес синего наименьший, поскольку колбочки, отвечающие за синий, малочисленны.

Если бы вместо взвешенной суммы использовалось простое среднее $(R+G+B)/3$, красный и зелёный цвета одинаковой интенсивности дали бы одинаковую яркость, что противоречит человеческому восприятию. Именно поэтому используются именно такие коэффициенты.

\section{Яркость, контраст, насыщенность}

Яркость изменяет все пиксели изображения на одну и ту же величину:
\[
Y' = Y + b,
\]
где $b \in [-120, 120]$ --- параметр, задаваемый ползунком. Положительные значения делают изображение светлее, отрицательные --- темнее. Если при этом значение выходит за пределы $[0, 255]$, оно отсекается (клиппинг): например, если $Y + b = 300$, то результат будет $255$. Это означает, что в светах теряется информация.

Контраст изменяет разброс значений относительно среднего:
\[
Y' = (Y - 128) \cdot c + 128,
\]
где $c > 0$ --- коэффициент контраста. При $c > 1$ разброс растягивается: тёмное становится темнее, светлое --- светлее. При $c < 1$ разброс сжимается, изображение становится блёклым. Если $c$ слишком велико, тени уходят в $0$, а света --- в $255$, и детали теряются.

Насыщенность изменяет интенсивность цвета, не затрагивая яркость. Для этого изображение переводится в пространство HSV, насыщенность умножается на коэффициент $s$, и результат переводится обратно в RGB:
\[
S' = S \cdot s.
\]
При $s = 0$ изображение становится полностью серым. При $s > 1$ цвета становятся более сочными, но если какая-то компонента упирается в максимум, оттенок может искажаться.

\section{Гистограмма}

Гистограмма изображения --- это функция $h(v)$, равная числу пикселей, яркость которых равна $v$. Формально, для изображения с $N$ пикселями и $256$ уровнями яркости:
\[
h(v) = \#\{ (x, y) : I(x, y) = v \}, \quad v = 0, \ldots, 255.
\]
Гистограмма показывает распределение яркостей. Узкий пик означает низкий контраст (все пиксели близки по яркости), широкое плато --- высокий контраст. Смещение влево означает, что изображение тёмное, смещение вправо --- светлое. Провалы в середине могут указывать на потерянные при сжатии уровни.

Для цветного изображения строятся три гистограммы --- по каналам $R$, $G$, $B$. Их сравнение позволяет понять цветовой баланс: если красный канал смещён вправо относительно синего, изображение имеет тёплый оттенок.

\section{Коррекция чёрно-белых изображений}

\subsection{Линейная коррекция}

Линейная коррекция растягивает диапазон яркостей $[b, w]$ на полный диапазон $[0, 255]$:
\[
Y' = \frac{Y - b}{w - b} \cdot 255,
\]
где $b$ --- чёрная точка, $w$ --- белая точка. Все значения меньше $b$ становятся чёрными, все значения больше $w$ --- белыми. Это полезно, если фотография выглядит серой из-за недостаточной экспозиции: правильный выбор $b$ и $w$ позволяет «вытянуть» контраст без потери деталей.

\subsection{Гамма-коррекция}

Гамма-коррекция --- нелинейное преобразование:
\[
Y' = 255 \left(\frac{Y}{255}\right)^{\gamma}.
\]
При $\gamma < 1$ тёмные области осветляются, при $\gamma > 1$ --- затемняются. Это соответствует изменению воспринимаемой яркости без изменения порядка значений: $\gamma$-коррекция не создаёт новых уровней, но перераспределяет существующие.

\subsection{Логарифмическая и экспоненциальная коррекции}

Логарифмическая коррекция:
\[
Y' = 255 \cdot \frac{\ln(1 + Y)}{\ln 256}
\]
сильно осветляет тёмные области и слабо влияет на светлые. Она применяется для сжатия динамического диапазона, когда нужно увидеть детали в тенях.

Экспоненциальная коррекция:
\[
Y' = 255 \cdot \frac{e^{kY/255} - 1}{e^k - 1}
\]
с параметром $k > 0$ действует наоборот: затемняет тёмные области, растягивая света. Применяется для усиления контраста в светах.

\section{Морфологические операции}

Морфология работает с формой объектов на изображении. Базовый инструмент --- \emph{структурирующий элемент}, небольшое бинарное ядро заданной формы (прямоугольник, крест, эллипс). Ядро скользит по изображению, и результат в каждой точке определяется взаимным расположением пикселей изображения и ядра.

\subsection{Эрозия}

Эрозия заменяет каждый пиксель на минимум в окне, заданном ядром:
\[
(I \ominus K)(x, y) = \min_{(i, j) \in K} I(x + i, y + j).
\]
Светлые объекты при этом «съедаются» по краям, тёмные расширяются. Эрозия убирает мелкие светлые детали и разделяет объекты, соединённые тонкими перемычками.

\subsection{Дилатация}

Дилатация заменяет каждый пиксель на максимум в окне:
\[
(I \oplus K)(x, y) = \max_{(i, j) \in K} I(x + i, y + j).
\]
Светлые объекты расширяются, тёмные сужаются. Дилатация заполняет мелкие тёмные трещины и соединяет близкие объекты.

\subsection{Открытие и закрытие}

Открытие --- это последовательное применение эрозии и дилатации:
\[
I \circ K = (I \ominus K) \oplus K.
\]
Открытие убирает мелкие светлые детали, но сохраняет общую форму крупных объектов. Оно используется для удаления шума в виде отдельных светлых точек.

Закрытие --- обратная операция:
\[
I \bullet K = (I \oplus K) \ominus K.
\]
Оно заполняет мелкие тёмные дырки и разрывы внутри светлых объектов.

\subsection{Градиент}

Морфологический градиент --- это разность между дилатацией и эрозией:
\[
\text{grad}(I) = (I \oplus K) - (I \ominus K).
\]
Он выделяет границы объектов: там, где яркость резко меняется, разность будет большой. Градиент даёт контуры толщиной, равной размеру ядра.

\section{Фильтры и свёртка}

Свёртка --- базовая операция линейной фильтрации. Результат в точке $(x, y)$ определяется как взвешенная сумма значений в окрестности:
\[
(I * K)(x, y) = \sum_{i=-m}^{m} \sum_{j=-n}^{n} I(x + i, y + j) \cdot K(i, j).
\]
Ядро $K$ задаёт, какие веса применяются к соседним пикселям.

\subsection{Повышение резкости}

Ядро для повышения резкости:
\[
K = \begin{pmatrix} 0 & -a & 0 \\ -a & 1 + 4a & -a \\ 0 & -a & 0 \end{pmatrix}.
\]
Центральный пиксель усиливается, соседние вычитаются. Чем больше параметр $a$, тем сильнее эффект. Резкость усиливается на границах, где соседние пиксели сильно отличаются.

\subsection{Размытие в движении}

Ядро представляет собой отрезок длиной $L$ под углом $\theta$:
\[
K(i, j) = \frac{1}{L} \text{ если } (i, j) \in \text{отрезок}, \quad 0 \text{ иначе}.
\]
Все пиксели вдоль отрезка усредняются с одинаковым весом, поэтому изображение «смазывается» в направлении движения.

\subsection{Тиснение}

Ядро тиснения:
\[
K = \begin{pmatrix} -2 & -1 & 0 \\ -1 & 1 & 1 \\ 0 & 1 & 2 \end{pmatrix}.
\]
Оно выделяет диагональные градиенты, создавая эффект выпуклости.

\subsection{Медианная фильтрация}

Медианный фильтр заменяет каждый пиксель на медиану значений в окне размером $k \times k$. Медиана устойчива к выбросам, поэтому фильтр хорошо убирает импульсный шум (одиночные белые или чёрные точки), не размывая края. В отличие от линейных фильтров, медианный фильтр нелинейный и не описывается свёрткой.

\section{Детекторы границ}

\subsection{Оператор Canny}

Детектор Canny --- многоступенчатый алгоритм. Сначала изображение сглаживается гауссовым фильтром для подавления шума. Затем вычисляется градиент по направлениям $x$ и $y$ через операторы Собеля, и находятся его магнитуда и направление:
\[
G = \sqrt{G_x^2 + G_y^2}, \quad \theta = \arctan\frac{G_y}{G_x}.
\]
Далее выполняется \emph{подавление немаксимумов}: пиксель оставляется только если его магнитуда больше, чем у соседей вдоль направления градиента. Наконец, применяется \emph{двойная пороговая фильтрация}: пиксели с магнитудой выше верхнего порога считаются «сильными» границами, между нижним и верхним --- «слабыми». Слабые сохраняются только если они соединены с сильными.

Пороги $T_1$ (нижний) и $T_2$ (верхний) задаются пользователем. Оптимальное соотношение --- $T_2 \approx 2 T_1$.

\subsection{Оператор Робертса}

Оператор Робертса использует две диагональные маски размером $2 \times 2$:
\[
K_x = \begin{pmatrix} 1 & 0 \\ 0 & -1 \end{pmatrix}, \quad
K_y = \begin{pmatrix} 0 & 1 \\ -1 & 0 \end{pmatrix}.
\]
Результат --- комбинация $\sqrt{(I * K_x)^2 + (I * K_y)^2}$. Оператор чувствителен к диагональным границам, но плохо подавляет шум, поэтому применяется редко.

\section{Ключевые точки}

Ключевые точки --- это характерные пиксели изображения, которые легко найти повторно при изменении масштаба, поворота или освещения. Они используются в задачах сопоставления изображений, панорамирования, отслеживания объектов.

\subsection{Детектор Харриса}

Детектор Харриса анализирует матрицу вторых моментов
\[
M = \sum_{(x, y) \in W} \begin{pmatrix} I_x^2 & I_x I_y \\ I_x I_y & I_y^2 \end{pmatrix},
\]
где $I_x$, $I_y$ --- производные изображения по $x$ и $y$, суммирование идёт по окну $W$ вокруг точки. Собственные числа $\lambda_1, \lambda_2$ матрицы $M$ показывают, насколько точка «угловатая»: если оба малы, точка лежит на плоской области; если одно мало, а другое велико --- на границе; если оба велики --- в углу. Вместо явного вычисления собственных чисел используется функция отклика
\[
R = \det M - k (\operatorname{tr} M)^2,
\]
где $k \approx 0{,}04$. Точки с большим $R$ объявляются углами.

\subsection{Детектор FAST}

FAST (Features from Accelerated Segment Test) проверяет окружность Брезенхэма радиуса 3 вокруг каждого пикселя (16 точек). Если среди них есть $n$ подряд идущих пикселей ярче или темнее центрального на порог $t$, точка считается угловой. Параметр $n$ обычно равен 9 или 12. FAST очень быстрый, но не инвариантен к масштабу.

\subsection{Детектор ORB}

ORB (Oriented FAST and Rotated BRIEF) объединяет FAST для поиска углов с вычислением ориентации через моменты изображения и построением вращательно-инвариантного дескриптора BRIEF. Он работает быстрее SIFT и часто используется в реальном времени.

\subsection{Детектор SIFT}

SIFT (Scale-Invariant Feature Transform) строит гауссову пирамиду: изображение последовательно размывается гауссовым фильтром с увеличивающейся сигмой $\sigma \cdot k^s$, где $k = 2^{1/3}$. Затем вычисляется разность гауссов (Difference of Gaussians, DoG):
\[
D(x, y, \sigma) = G(x, y, k\sigma) - G(x, y, \sigma),
\]
которая аппроксимирует нормализованный лапласиан гауссиана. Экстремумы в трёхмерном пространстве $(x, y, \sigma)$ --- кандидаты в ключевые точки. Каждая точка сравнивается с $26$ соседями: $9$ на предыдущем масштабе, $8$ на текущем, $9$ на следующем. Если она больше или меньше всех --- она ключевая. Затем слабые точки отбрасываются по порогу контраста, а точки на границах --- по соотношению главных кривизн. SIFT инвариантен к масштабу и вращению, что делает его одним из лучших детекторов для сопоставления изображений.

\section{Поиск похожих изображений}

Для сравнения изображений используется гистограмма в пространстве HSV. Каждое изображение описывается двумерной гистограммой $H(h, s)$ размером $30 \times 32$ корзины по оттенку и насыщенности. Гистограммы нормализуются так, чтобы сумма значений равнялась единице.

Мера похожести --- коэффициент корреляции Пирсона между двумя гистограммами $H_1$ и $H_2$:
\[
\rho(H_1, H_2) = \frac{\sum_i (H_1(i) - \bar{H}_1)(H_2(i) - \bar{H}_2)}{\sqrt{\sum_i (H_1(i) - \bar{H}_1)^2} \sqrt{\sum_i (H_2(i) - \bar{H}_2)^2}}.
\]
Значение $\rho = 1$ означает полное совпадение, $\rho = 0$ --- отсутствие корреляции. Из всех пар изображений выбирается пара с максимальным $\rho$.

\section{Обработка видео}

\subsection{Вычитание фона}

Вычитание фона --- метод выделения движущихся объектов на видео со статичной камерой. Для каждого пикселя строится статистическая модель фона на основе предыдущих кадров. Затем для каждого нового кадра вычисляется, насколько пиксель отличается от модели; если отличие превышает порог, пиксель считается частью переднего плана.

Алгоритм MOG2 (Mixture of Gaussians) моделирует распределение значений каждого пикселя как смесь нескольких гауссовых распределений:
\[
P(I) = \sum_{k=1}^{K} \omega_k \cdot \mathcal{N}(I \mid \mu_k, \sigma_k^2),
\]
где $\omega_k$ --- вес $k$-й компоненты, $\mu_k$ и $\sigma_k$ --- её среднее и стандартное отклонение. Пиксель считается фоновым, если его значение попадает в одну из «фоновых» компонент с большим весом. Модель постепенно обновляется, поэтому медленные изменения освещения не вызывают ложных срабатываний.

Алгоритм KNN использует не смесь гауссиан, а набор из $K$ ближайших образцов цвета для каждого пикселя. Он устойчивее к шуму и лучше работает при сложном фоне.

Результат вычитания фона --- бинарная маска, в которой белые пиксели принадлежат движущимся объектам, а чёрные --- фону. Умножая оригинальный кадр на маску, получают изображение, на котором движущиеся объекты сохраняют цвет, а фон становится чёрным. Это и есть режим «Только объекты на чёрном фоне».

\subsection{Размытие движущихся объектов}

Размытие движущихся объектов использует ту же маску переднего плана. Оригинальный кадр размывается гауссовым фильтром; затем к размытому изображению применяется маска (движущиеся объекты становятся размытыми), а к оригиналу --- инвертированная маска (фон остаётся резким). Результат:
\[
\text{out}(x, y) = \begin{cases}
\text{blur}(x, y), & \text{если маска}(x, y) = 255, \\
\text{orig}(x, y), & \text{иначе}.
\end{cases}
\]
Эффект подчёркивает движущиеся объекты, привлекая к ним внимание.

\section{Заключение}

Все рассмотренные операции можно разделить на четыре группы. Первая группа --- \emph{поточечные} операции: яркость, контраст, гамма-коррекция, преобразование в градации серого. Здесь значение нового пикселя зависит только от старого значения. Вторая группа --- \emph{локальные} операции: свёртки, медианный фильтр, морфология. Значение зависит от окрестности пикселя. Третья группа --- \emph{глобальные} операции: гистограмма, коррекция по чёрной и белой точкам. Здесь новое значение зависит от свойств всего изображения. Четвёртая группа --- \emph{структурные} операции: детекторы границ, ключевые точки, поиск похожих изображений. Они не изменяют изображение, а выделяют из него информацию.

Понимание этой иерархии позволяет осознанно выбирать инструменты: для подавления шума --- локальные фильтры, для улучшения видимости --- поточечные и глобальные операции, для анализа содержимого --- структурные методы. Навык выбирать правильный инструмент под задачу приходит с практикой; математическая основа всех операций изложена выше и служит справочником при работе с любым программным обеспечением для обработки изображений.

\end{document}